\documentclass{umk-matematika}

\begin{document}

\maketitlepage
    {Практическая работа №24}
    {Логарифмические уравнения}
    {}

\section{Фундамент (1--4 балла)}

\textbf{Задача 1 (1 балл).} Решите уравнение: $\log_{2} x = 3$

\textbf{Задача 2 (2 балла).} Решите уравнение: $\log_{5} (x - 1) = 2$

\textbf{Задача 3 (3 балла).} Решите уравнение: $\log_{3} (3x) = 2$

\textbf{Задача 4 (4 балла).} Решите уравнение: $\log_{4} (x + 3) = 1$

\section{Стены (5--7 баллов)}

\textbf{Задача 5 (5 баллов).} Решите уравнение: $\log_{2} (x + 1) + \log_{2} (x - 1) = 3$

\textbf{Задача 6 (6 баллов).} Решите уравнение: $\log_{3} (x^2 - 4) - \log_{3} (x - 2) = 1$

\textbf{Задача 7 (7 баллов).} Решите уравнение: $\log_{2}^{2} x - 3\log_{2} x + 2 = 0$

\section{Кровля (8--10 баллов)}

\textbf{Задача 8 (8 баллов).} Срок службы оборудования $t$ (лет) связан с износом $I$ (\%) формулой $t = 5 \cdot \log_{2} \left(\frac{100}{100 - I}\right)$. Через сколько лет износ достигнет $75\%$?

\textbf{Задача 9 (9 баллов).} Решите уравнение: $\log_{2} (x + 2) + \log_{2} (x - 2) = \log_{2} (5x - 10)$

\textbf{Задача 10 (10 баллов).} Решите уравнение: $\log_{5} x + \log_{x} 5 = 2$

\newpage
\section*{Ответы}

\begin{enumerate}
  \item $x = 8$
  \item $x = 26$
  \item $x = 3$
  \item $x = 1$
  \item $x = 3$
  \item нет решений
  \item $x = 2;\ x = 4$
  \item через $10$ лет
  \item $x = 3$
  \item $x = 5$
\end{enumerate}

\end{document}